% =============================================================================
% One-pager A4 - G1: LangGraph e Orquestracao Baseada em Grafos
% Para distribuir aos visitantes da sessao de paineis.
% Compilar: pdflatex one_pager.tex (no Overleaf, defina como main document)
% Imagens em metrics/ (mesmo esquema do painel).
% =============================================================================

\documentclass[a4paper,10pt]{article}

\usepackage[utf8]{inputenc}
\usepackage[T1]{fontenc}
\usepackage[brazil]{babel}
\usepackage[a4paper, margin=12mm]{geometry}
\usepackage{graphicx}
\usepackage{booktabs}
\usepackage{multicol}
\usepackage{xcolor}
\usepackage{titlesec}
\usepackage{enumitem}
\usepackage{hyperref}
\usepackage{lmodern}

\definecolor{lgPrimary}{HTML}{1F3A93}
\definecolor{lgSoft}{HTML}{F4F6FB}

\hypersetup{colorlinks=true, urlcolor=lgPrimary, linkcolor=lgPrimary}

\setlist{nosep, leftmargin=*}
\titleformat{\section}{\color{lgPrimary}\bfseries\large}{}{0pt}{}
\titleformat{\subsection}{\color{lgPrimary}\bfseries\normalsize}{}{0pt}{}
\titlespacing*{\section}{0pt}{6pt}{2pt}
\titlespacing*{\subsection}{0pt}{4pt}{2pt}

\setlength{\columnsep}{6mm}
\setlength{\parindent}{0pt}
\setlength{\parskip}{2pt}

\pagestyle{empty}

\begin{document}

% =============================================================================
% FRENTE
% =============================================================================
\begin{center}
{\Large\bfseries\color{lgPrimary} LangGraph e Orquestração Baseada em Grafos}\\[2pt]
{\small Agentes de IA -- Grupo G1 \quad $|$ \quad Cecilia Rufatto, Felipe Proença, Leonardo Moino, Leonardo Quartaroli, Lorena Scabello}
\end{center}

\vspace{2pt}\hrule\vspace{4pt}

\begin{multicols}{2}

\section*{Contexto}
Frameworks lineares como o LangChain tradicional executam etapas em ordem fixa, sem suporte nativo a ciclos, estado persistente ou decisões condicionais. O LangGraph cobre essa lacuna tratando o agente como um grafo de estados: nós são funções sobre o estado, arestas decidem o caminho e \emph{checkpointers} oficiais cuidam da persistência.

\section*{Conceitos-chave}
\begin{itemize}
\item \textbf{StateGraph} --- grafo cujo estado é um \texttt{TypedDict} compartilhado entre nós.
\item \textbf{Node} --- função Python que recebe e devolve o estado.
\item \textbf{Conditional Edge} --- roteamento dinâmico via \texttt{router(state)}.
\item \textbf{Reducer} --- estratégia de \emph{merge} para campos do estado.
\item \textbf{Checkpointer} --- persistência entre execuções (SQLite, Postgres, Redis).
\item \textbf{HITL} --- pausa via \texttt{interrupt} e retoma com \texttt{Command}.
\end{itemize}

\columnbreak

\section*{Arquitetura do protótipo}
\centering
\includegraphics[width=0.95\linewidth]{metrics/graph.png}\\[2pt]
\raggedright\small
Cada mensagem é processada pelo \texttt{detect\_intent}, rotulada como \emph{pedido}, \emph{suporte} ou \emph{geral}, e despachada pela aresta condicional \texttt{route\_intent} pro handler. Cada handler grava o histórico em JSON pelo \texttt{memory.py}.

\end{multicols}

\vspace{2pt}\hrule\vspace{4pt}

\section*{Resultados em um dataset de 40 mensagens (pt-BR)}

\begin{multicols}{2}

\begin{tabular}{@{} l r @{}}
\toprule
\textbf{Métrica}      & \textbf{Valor} \\
\midrule
Accuracy              & \textbf{85,00\%} \\
Macro F1              & 0,8507 \\
Latência média        & 33,52 ms \\
Latência P95          & $\sim$56 ms \\
\bottomrule
\end{tabular}

\vspace{4pt}
\begin{tabular}{@{} l r r r @{}}
\toprule
\textbf{Classe} & \textbf{P} & \textbf{R} & \textbf{F1} \\
\midrule
pedido   & 0,92 & 0,80 & 0,86 \\
suporte  & 0,92 & 0,85 & 0,88 \\
geral    & 0,73 & 0,92 & 0,81 \\
\bottomrule
\end{tabular}

\columnbreak

\centering
\includegraphics[width=0.85\linewidth]{metrics/confusion_matrix.png}

\end{multicols}

\vfill
\begin{center}
\small\color{lgPrimary}\textbf{Verso:} estado da arte, limitações, próximos passos e referências.
\end{center}

\newpage

% =============================================================================
% VERSO
% =============================================================================
\begin{center}
{\Large\bfseries\color{lgPrimary} Estado da arte, discussão e referências}
\end{center}

\vspace{2pt}\hrule\vspace{4pt}

\section*{Estado da arte}

\begin{multicols}{2}

\subsection*{Linha do tempo}
\begin{itemize}
\item \textbf{2022} -- LangChain lança \emph{chains} lineares.
\item \textbf{Jan/2024} -- LangGraph 0.0.x: primeira versão pública.
\item \textbf{Mai/2024} -- LangGraph Studio + Cloud.
\item \textbf{2025} -- LangGraph 1.x estável com \emph{durable execution}.
\end{itemize}

\columnbreak

\subsection*{Comparativo}
\begin{tabular}{@{} p{0.30\linewidth} p{0.30\linewidth} p{0.30\linewidth} @{}}
\toprule
\textbf{Framework} & \textbf{Ciclos} & \textbf{Persistência} \\
\midrule
LangGraph    & Nativo   & Checkpointers \\
CrewAI       & Limitado & Memory mod. \\
AutoGen      & Por chat & Externa \\
Swarm        & Não      & Não \\
\bottomrule
\end{tabular}

\end{multicols}

\vspace{2pt}\hrule\vspace{4pt}

\section*{O que funcionou}
A combinação de StateGraph com aresta condicional deixou o roteamento declarativo e fácil de testar --- dá pra escrever evals direto sobre o grafo compilado. Persistência em JSON foi suficiente para o protótipo, e a latência abaixo de 60 ms no P95 mantém o agente confortável em chat síncrono.

\section*{Limitações observadas}
A classificação por keywords não captura semântica: a classe \texttt{geral} acaba absorvendo qualquer frase que não bate com as outras --- daí o \emph{precision} de 0,73. Mensagens multi-intenção, como ``meu produto chegou quebrado e quero suporte'', são forçadas a uma única rota. Também não tem suporte a contexto profundo entre turnos.

\section*{Próximos passos}
Trocar o \texttt{detect\_intent} por um classificador LLM \emph{few-shot} com saída estruturada via Pydantic; adicionar um nó de Human-in-the-Loop para casos em que o \texttt{confidence} fique baixo; e migrar a persistência JSON para um \texttt{SqliteSaver}, ganhando \emph{durable execution}.

\vspace{2pt}\hrule\vspace{4pt}

\section*{Referências principais (ABNT NBR 6023)}
\footnotesize
\begin{enumerate}
\item LANGCHAIN AI. \textbf{LangGraph -- Overview and Graph API}. Documentação oficial, 2024. \texttt{langchain-ai.github.io/langgraph/}.
\item LANGCHAIN AI. \textbf{Persistence and Checkpointers}. 2024. \texttt{langchain-ai.github.io/langgraph/concepts/persistence/}.
\item LANGCHAIN AI. \textbf{Human-in-the-loop}. 2024. \texttt{langchain-ai.github.io/langgraph/concepts/human\_in\_the\_loop/}.
\item WU, Q. \emph{et al.} \textbf{AutoGen: Enabling Next-Gen LLM Applications via Multi-Agent Conversation}. arXiv:2308.08155, 2023.
\item YAO, S. \emph{et al.} \textbf{ReAct: Synergizing Reasoning and Acting in Language Models}. ICLR, 2023.
\item OPENAI. \textbf{Swarm}. GitHub, 2024.
\item LANGCHAIN AI. \textbf{LangGraph Academy}. \texttt{academy.langchain.com/courses/intro-to-langgraph}.
\end{enumerate}

\normalsize
\vspace{4pt}\hrule\vspace{4pt}

\begin{minipage}[c]{0.70\linewidth}
\textbf{\color{lgPrimary} Repositório do grupo}\\
\texttt{github.com/Ceciliarufatto/langgraph}\\[2pt]
Código-fonte, dataset, métricas, painel e demo em GIF.
\end{minipage}\hfill
\begin{minipage}[c]{0.25\linewidth}
\centering

\includegraphics[width=0.85\linewidth]{metrics/qrcode_repo.png}
\end{minipage}

\end{document}
